\documentclass[11pt]{amsart}
\usepackage[margin=1.1in]{geometry}
\usepackage{amssymb,mathtools,booktabs,multicol}
\usepackage[hidelinks]{hyperref}
\newcommand{\nclosed}{243}
\newcommand{\nopen}{293}
\newcommand{\ntotal}{536}
\newcommand{\grouptable}{$\O_4(\F_{5})$ & reflections & 15 \\
$\Sp_4(\F_{3})$ & transvections & 67 \\
$\Sp_4(\F_{5})$ & transvections & 161 \\}
\newcommand{\closedlist}{$13\mathrm{a}_{69}$, $13\mathrm{a}_{70}$, $13\mathrm{a}_{72}$, $13\mathrm{a}_{73}$, $13\mathrm{a}_{74}$, $13\mathrm{a}_{76}$, $13\mathrm{a}_{79}$, $13\mathrm{a}_{96}$, $13\mathrm{a}_{104}$, $13\mathrm{a}_{106}$, $13\mathrm{a}_{107}$, $13\mathrm{a}_{108}$, $13\mathrm{a}_{156}$, $13\mathrm{a}_{157}$, $13\mathrm{a}_{160}$, $13\mathrm{a}_{171}$, $13\mathrm{a}_{235}$, $13\mathrm{a}_{343}$, $13\mathrm{a}_{353}$, $13\mathrm{a}_{354}$, $13\mathrm{a}_{355}$, $13\mathrm{a}_{358}$, $13\mathrm{a}_{359}$, $13\mathrm{a}_{360}$, $13\mathrm{a}_{388}$, $13\mathrm{a}_{390}$, $13\mathrm{a}_{396}$, $13\mathrm{a}_{399}$, $13\mathrm{a}_{400}$, $13\mathrm{a}_{402}$, $13\mathrm{a}_{403}$, $13\mathrm{a}_{408}$, $13\mathrm{a}_{413}$, $13\mathrm{a}_{414}$, $13\mathrm{a}_{505}$, $13\mathrm{a}_{524}$, $13\mathrm{a}_{525}$, $13\mathrm{a}_{526}$, $13\mathrm{a}_{527}$, $13\mathrm{a}_{528}$, $13\mathrm{a}_{529}$, $13\mathrm{a}_{535}$, $13\mathrm{a}_{536}$, $13\mathrm{a}_{541}$, $13\mathrm{a}_{542}$, $13\mathrm{a}_{543}$, $13\mathrm{a}_{544}$, $13\mathrm{a}_{545}$, $13\mathrm{a}_{554}$, $13\mathrm{a}_{555}$, $13\mathrm{a}_{557}$, $13\mathrm{a}_{558}$, $13\mathrm{a}_{596}$, $13\mathrm{a}_{597}$, $13\mathrm{a}_{601}$, $13\mathrm{a}_{621}$, $13\mathrm{a}_{630}$, $13\mathrm{a}_{631}$, $13\mathrm{a}_{632}$, $13\mathrm{a}_{636}$, $13\mathrm{a}_{637}$, $13\mathrm{a}_{657}$, $13\mathrm{a}_{664}$, $13\mathrm{a}_{665}$, $13\mathrm{a}_{815}$, $13\mathrm{a}_{867}$, $13\mathrm{a}_{1392}$, $13\mathrm{a}_{1538}$, $13\mathrm{a}_{1995}$, $13\mathrm{a}_{2006}$, $13\mathrm{a}_{2008}$, $13\mathrm{a}_{2613}$, $13\mathrm{a}_{2614}$, $13\mathrm{a}_{2615}$, $13\mathrm{a}_{2616}$, $13\mathrm{a}_{2617}$, $13\mathrm{a}_{2619}$, $13\mathrm{a}_{2620}$, $13\mathrm{a}_{2800}$, $13\mathrm{a}_{2801}$, $13\mathrm{a}_{3038}$, $13\mathrm{a}_{3040}$, $13\mathrm{a}_{3107}$, $13\mathrm{a}_{3144}$, $13\mathrm{a}_{3146}$, $13\mathrm{a}_{3179}$, $13\mathrm{a}_{3257}$, $13\mathrm{a}_{3333}$, $13\mathrm{a}_{3334}$, $13\mathrm{a}_{3343}$, $13\mathrm{a}_{3344}$, $13\mathrm{a}_{3376}$, $13\mathrm{a}_{3377}$, $13\mathrm{a}_{3378}$, $13\mathrm{a}_{3380}$, $13\mathrm{a}_{3381}$, $13\mathrm{n}_{160}$, $13\mathrm{n}_{161}$, $13\mathrm{n}_{162}$, $13\mathrm{n}_{163}$, $13\mathrm{n}_{164}$, $13\mathrm{n}_{165}$, $13\mathrm{n}_{166}$, $13\mathrm{n}_{167}$, $13\mathrm{n}_{168}$, $13\mathrm{n}_{169}$, $13\mathrm{n}_{173}$, $13\mathrm{n}_{174}$, $13\mathrm{n}_{175}$, $13\mathrm{n}_{176}$, $13\mathrm{n}_{177}$, $13\mathrm{n}_{178}$, $13\mathrm{n}_{179}$, $13\mathrm{n}_{180}$, $13\mathrm{n}_{181}$, $13\mathrm{n}_{182}$, $13\mathrm{n}_{211}$, $13\mathrm{n}_{429}$, $13\mathrm{n}_{430}$, $13\mathrm{n}_{431}$, $13\mathrm{n}_{432}$, $13\mathrm{n}_{433}$, $13\mathrm{n}_{434}$, $13\mathrm{n}_{435}$, $13\mathrm{n}_{440}$, $13\mathrm{n}_{441}$, $13\mathrm{n}_{442}$, $13\mathrm{n}_{449}$, $13\mathrm{n}_{467}$, $13\mathrm{n}_{601}$, $13\mathrm{n}_{609}$, $13\mathrm{n}_{610}$, $13\mathrm{n}_{929}$, $13\mathrm{n}_{930}$, $13\mathrm{n}_{931}$, $13\mathrm{n}_{932}$, $13\mathrm{n}_{933}$, $13\mathrm{n}_{957}$, $13\mathrm{n}_{958}$, $13\mathrm{n}_{959}$, $13\mathrm{n}_{960}$, $13\mathrm{n}_{961}$, $13\mathrm{n}_{962}$, $13\mathrm{n}_{963}$, $13\mathrm{n}_{964}$, $13\mathrm{n}_{965}$, $13\mathrm{n}_{966}$, $13\mathrm{n}_{967}$, $13\mathrm{n}_{968}$, $13\mathrm{n}_{977}$, $13\mathrm{n}_{978}$, $13\mathrm{n}_{979}$, $13\mathrm{n}_{980}$, $13\mathrm{n}_{981}$, $13\mathrm{n}_{982}$, $13\mathrm{n}_{983}$, $13\mathrm{n}_{1042}$, $13\mathrm{n}_{1058}$, $13\mathrm{n}_{1059}$, $13\mathrm{n}_{1060}$, $13\mathrm{n}_{1061}$, $13\mathrm{n}_{1062}$, $13\mathrm{n}_{1063}$, $13\mathrm{n}_{1064}$, $13\mathrm{n}_{1072}$, $13\mathrm{n}_{1073}$, $13\mathrm{n}_{1074}$, $13\mathrm{n}_{1075}$, $13\mathrm{n}_{1076}$, $13\mathrm{n}_{1077}$, $13\mathrm{n}_{1078}$, $13\mathrm{n}_{1082}$, $13\mathrm{n}_{1083}$, $13\mathrm{n}_{1084}$, $13\mathrm{n}_{1103}$, $13\mathrm{n}_{1104}$, $13\mathrm{n}_{1105}$, $13\mathrm{n}_{1106}$, $13\mathrm{n}_{1239}$, $13\mathrm{n}_{1280}$, $13\mathrm{n}_{1281}$, $13\mathrm{n}_{1282}$, $13\mathrm{n}_{1283}$, $13\mathrm{n}_{1284}$, $13\mathrm{n}_{1285}$, $13\mathrm{n}_{1286}$, $13\mathrm{n}_{1287}$, $13\mathrm{n}_{1288}$, $13\mathrm{n}_{1289}$, $13\mathrm{n}_{1290}$, $13\mathrm{n}_{1291}$, $13\mathrm{n}_{1292}$, $13\mathrm{n}_{1293}$, $13\mathrm{n}_{1294}$, $13\mathrm{n}_{1295}$, $13\mathrm{n}_{1305}$, $13\mathrm{n}_{1306}$, $13\mathrm{n}_{1307}$, $13\mathrm{n}_{1308}$, $13\mathrm{n}_{1309}$, $13\mathrm{n}_{1310}$, $13\mathrm{n}_{1311}$, $13\mathrm{n}_{1312}$, $13\mathrm{n}_{1313}$, $13\mathrm{n}_{1314}$, $13\mathrm{n}_{1315}$, $13\mathrm{n}_{1316}$, $13\mathrm{n}_{1317}$, $13\mathrm{n}_{1318}$, $13\mathrm{n}_{1319}$, $13\mathrm{n}_{1320}$, $13\mathrm{n}_{1321}$, $13\mathrm{n}_{1322}$, $13\mathrm{n}_{1323}$, $13\mathrm{n}_{1324}$, $13\mathrm{n}_{1332}$, $13\mathrm{n}_{1333}$, $13\mathrm{n}_{1334}$, $13\mathrm{n}_{1363}$, $13\mathrm{n}_{1377}$, $13\mathrm{n}_{1378}$, $13\mathrm{n}_{1379}$, $13\mathrm{n}_{1380}$, $13\mathrm{n}_{1381}$, $13\mathrm{n}_{1385}$, $13\mathrm{n}_{1386}$, $13\mathrm{n}_{2933}$, $13\mathrm{n}_{2934}$, $13\mathrm{n}_{2935}$, $13\mathrm{n}_{2936}$, $13\mathrm{n}_{2937}$, $13\mathrm{n}_{2954}$, $13\mathrm{n}_{2955}$, $13\mathrm{n}_{2956}$, $13\mathrm{n}_{3390}$, $13\mathrm{n}_{3505}$, $13\mathrm{n}_{3723}$, $13\mathrm{n}_{3725}$, $13\mathrm{n}_{3726}$, $13\mathrm{n}_{3727}$, $13\mathrm{n}_{3763}$, $13\mathrm{n}_{3838}$, $13\mathrm{n}_{3915}$}
\newcommand{\openlist}{$13\mathrm{a}_{129}$, $13\mathrm{a}_{135}$, $13\mathrm{a}_{136}$, $13\mathrm{a}_{137}$, $13\mathrm{a}_{194}$, $13\mathrm{a}_{196}$, $13\mathrm{a}_{197}$, $13\mathrm{a}_{198}$, $13\mathrm{a}_{233}$, $13\mathrm{a}_{234}$, $13\mathrm{a}_{239}$, $13\mathrm{a}_{242}$, $13\mathrm{a}_{243}$, $13\mathrm{a}_{244}$, $13\mathrm{a}_{253}$, $13\mathrm{a}_{265}$, $13\mathrm{a}_{268}$, $13\mathrm{a}_{269}$, $13\mathrm{a}_{270}$, $13\mathrm{a}_{271}$, $13\mathrm{a}_{272}$, $13\mathrm{a}_{304}$, $13\mathrm{a}_{309}$, $13\mathrm{a}_{338}$, $13\mathrm{a}_{362}$, $13\mathrm{a}_{365}$, $13\mathrm{a}_{366}$, $13\mathrm{a}_{369}$, $13\mathrm{a}_{373}$, $13\mathrm{a}_{375}$, $13\mathrm{a}_{387}$, $13\mathrm{a}_{391}$, $13\mathrm{a}_{392}$, $13\mathrm{a}_{471}$, $13\mathrm{a}_{483}$, $13\mathrm{a}_{487}$, $13\mathrm{a}_{516}$, $13\mathrm{a}_{518}$, $13\mathrm{a}_{530}$, $13\mathrm{a}_{623}$, $13\mathrm{a}_{629}$, $13\mathrm{a}_{3145}$, $13\mathrm{a}_{3147}$, $13\mathrm{a}_{3148}$, $13\mathrm{a}_{3149}$, $13\mathrm{a}_{3167}$, $13\mathrm{a}_{3168}$, $13\mathrm{a}_{3169}$, $13\mathrm{a}_{3170}$, $13\mathrm{a}_{3188}$, $13\mathrm{a}_{3189}$, $13\mathrm{a}_{3190}$, $13\mathrm{a}_{3191}$, $13\mathrm{a}_{3204}$, $13\mathrm{a}_{3205}$, $13\mathrm{a}_{3207}$, $13\mathrm{a}_{3208}$, $13\mathrm{a}_{3209}$, $13\mathrm{a}_{3210}$, $13\mathrm{a}_{3213}$, $13\mathrm{a}_{3216}$, $13\mathrm{a}_{3219}$, $13\mathrm{a}_{3220}$, $13\mathrm{a}_{3221}$, $13\mathrm{a}_{3222}$, $13\mathrm{a}_{3223}$, $13\mathrm{a}_{3232}$, $13\mathrm{a}_{3234}$, $13\mathrm{a}_{3248}$, $13\mathrm{a}_{3259}$, $13\mathrm{a}_{3261}$, $13\mathrm{a}_{3262}$, $13\mathrm{a}_{3263}$, $13\mathrm{a}_{3264}$, $13\mathrm{a}_{3266}$, $13\mathrm{a}_{3274}$, $13\mathrm{a}_{3275}$, $13\mathrm{a}_{3276}$, $13\mathrm{a}_{3291}$, $13\mathrm{a}_{3292}$, $13\mathrm{a}_{3293}$, $13\mathrm{a}_{3294}$, $13\mathrm{a}_{3295}$, $13\mathrm{a}_{3298}$, $13\mathrm{a}_{3301}$, $13\mathrm{a}_{3302}$, $13\mathrm{a}_{3303}$, $13\mathrm{a}_{3317}$, $13\mathrm{a}_{3320}$, $13\mathrm{a}_{3323}$, $13\mathrm{a}_{3330}$, $13\mathrm{a}_{3331}$, $13\mathrm{a}_{3335}$, $13\mathrm{a}_{3345}$, $13\mathrm{a}_{3346}$, $13\mathrm{a}_{3347}$, $13\mathrm{a}_{3348}$, $13\mathrm{a}_{3351}$, $13\mathrm{a}_{3354}$, $13\mathrm{a}_{3355}$, $13\mathrm{a}_{3356}$, $13\mathrm{a}_{3363}$, $13\mathrm{a}_{3365}$, $13\mathrm{a}_{3368}$, $13\mathrm{a}_{3374}$, $13\mathrm{a}_{3375}$, $13\mathrm{a}_{3388}$, $13\mathrm{a}_{3390}$, $13\mathrm{a}_{3393}$, $13\mathrm{a}_{3394}$, $13\mathrm{a}_{3395}$, $13\mathrm{a}_{3396}$, $13\mathrm{a}_{3399}$, $13\mathrm{a}_{3400}$, $13\mathrm{a}_{3401}$, $13\mathrm{a}_{3402}$, $13\mathrm{a}_{3403}$, $13\mathrm{a}_{3404}$, $13\mathrm{a}_{3405}$, $13\mathrm{a}_{3407}$, $13\mathrm{a}_{3409}$, $13\mathrm{a}_{3410}$, $13\mathrm{a}_{3411}$, $13\mathrm{a}_{3412}$, $13\mathrm{a}_{3413}$, $13\mathrm{a}_{3415}$, $13\mathrm{a}_{3417}$, $13\mathrm{a}_{3418}$, $13\mathrm{a}_{3811}$, $13\mathrm{a}_{3812}$, $13\mathrm{n}_{326}$, $13\mathrm{n}_{329}$, $13\mathrm{n}_{330}$, $13\mathrm{n}_{341}$, $13\mathrm{n}_{344}$, $13\mathrm{n}_{345}$, $13\mathrm{n}_{348}$, $13\mathrm{n}_{351}$, $13\mathrm{n}_{352}$, $13\mathrm{n}_{355}$, $13\mathrm{n}_{357}$, $13\mathrm{n}_{358}$, $13\mathrm{n}_{510}$, $13\mathrm{n}_{513}$, $13\mathrm{n}_{514}$, $13\mathrm{n}_{518}$, $13\mathrm{n}_{520}$, $13\mathrm{n}_{521}$, $13\mathrm{n}_{524}$, $13\mathrm{n}_{525}$, $13\mathrm{n}_{526}$, $13\mathrm{n}_{529}$, $13\mathrm{n}_{530}$, $13\mathrm{n}_{531}$, $13\mathrm{n}_{582}$, $13\mathrm{n}_{590}$, $13\mathrm{n}_{591}$, $13\mathrm{n}_{598}$, $13\mathrm{n}_{599}$, $13\mathrm{n}_{600}$, $13\mathrm{n}_{627}$, $13\mathrm{n}_{632}$, $13\mathrm{n}_{633}$, $13\mathrm{n}_{644}$, $13\mathrm{n}_{647}$, $13\mathrm{n}_{648}$, $13\mathrm{n}_{651}$, $13\mathrm{n}_{652}$, $13\mathrm{n}_{661}$, $13\mathrm{n}_{662}$, $13\mathrm{n}_{663}$, $13\mathrm{n}_{694}$, $13\mathrm{n}_{695}$, $13\mathrm{n}_{696}$, $13\mathrm{n}_{704}$, $13\mathrm{n}_{705}$, $13\mathrm{n}_{706}$, $13\mathrm{n}_{707}$, $13\mathrm{n}_{708}$, $13\mathrm{n}_{709}$, $13\mathrm{n}_{712}$, $13\mathrm{n}_{713}$, $13\mathrm{n}_{716}$, $13\mathrm{n}_{717}$, $13\mathrm{n}_{718}$, $13\mathrm{n}_{834}$, $13\mathrm{n}_{838}$, $13\mathrm{n}_{839}$, $13\mathrm{n}_{857}$, $13\mathrm{n}_{926}$, $13\mathrm{n}_{991}$, $13\mathrm{n}_{992}$, $13\mathrm{n}_{993}$, $13\mathrm{n}_{1000}$, $13\mathrm{n}_{1001}$, $13\mathrm{n}_{1004}$, $13\mathrm{n}_{1005}$, $13\mathrm{n}_{1006}$, $13\mathrm{n}_{1007}$, $13\mathrm{n}_{1011}$, $13\mathrm{n}_{1012}$, $13\mathrm{n}_{1013}$, $13\mathrm{n}_{1039}$, $13\mathrm{n}_{1040}$, $13\mathrm{n}_{1041}$, $13\mathrm{n}_{1043}$, $13\mathrm{n}_{1044}$, $13\mathrm{n}_{1045}$, $13\mathrm{n}_{1046}$, $13\mathrm{n}_{1047}$, $13\mathrm{n}_{1048}$, $13\mathrm{n}_{1175}$, $13\mathrm{n}_{1178}$, $13\mathrm{n}_{1179}$, $13\mathrm{n}_{1200}$, $13\mathrm{n}_{1201}$, $13\mathrm{n}_{1202}$, $13\mathrm{n}_{1210}$, $13\mathrm{n}_{1211}$, $13\mathrm{n}_{1212}$, $13\mathrm{n}_{1270}$, $13\mathrm{n}_{1276}$, $13\mathrm{n}_{1296}$, $13\mathrm{n}_{1367}$, $13\mathrm{n}_{1370}$, $13\mathrm{n}_{1371}$, $13\mathrm{n}_{1374}$, $13\mathrm{n}_{1375}$, $13\mathrm{n}_{1376}$, $13\mathrm{n}_{3739}$, $13\mathrm{n}_{3742}$, $13\mathrm{n}_{3745}$, $13\mathrm{n}_{3748}$, $13\mathrm{n}_{3770}$, $13\mathrm{n}_{3773}$, $13\mathrm{n}_{3776}$, $13\mathrm{n}_{3779}$, $13\mathrm{n}_{3814}$, $13\mathrm{n}_{3840}$, $13\mathrm{n}_{3842}$, $13\mathrm{n}_{3844}$, $13\mathrm{n}_{3854}$, $13\mathrm{n}_{3855}$, $13\mathrm{n}_{3856}$, $13\mathrm{n}_{3870}$, $13\mathrm{n}_{3871}$, $13\mathrm{n}_{3872}$, $13\mathrm{n}_{3873}$, $13\mathrm{n}_{3874}$, $13\mathrm{n}_{3897}$, $13\mathrm{n}_{3898}$, $13\mathrm{n}_{3899}$, $13\mathrm{n}_{3900}$, $13\mathrm{n}_{3901}$, $13\mathrm{n}_{3903}$, $13\mathrm{n}_{3904}$, $13\mathrm{n}_{3905}$, $13\mathrm{n}_{3911}$, $13\mathrm{n}_{3914}$, $13\mathrm{n}_{3916}$, $13\mathrm{n}_{3918}$, $13\mathrm{n}_{3919}$, $13\mathrm{n}_{3920}$, $13\mathrm{n}_{3921}$, $13\mathrm{n}_{3922}$, $13\mathrm{n}_{3924}$, $13\mathrm{n}_{3925}$, $13\mathrm{n}_{3928}$, $13\mathrm{n}_{3929}$, $13\mathrm{n}_{3931}$, $13\mathrm{n}_{3932}$, $13\mathrm{n}_{3933}$, $13\mathrm{n}_{3934}$, $13\mathrm{n}_{3935}$, $13\mathrm{n}_{3936}$, $13\mathrm{n}_{3937}$, $13\mathrm{n}_{3938}$, $13\mathrm{n}_{3939}$, $13\mathrm{n}_{3940}$, $13\mathrm{n}_{3941}$, $13\mathrm{n}_{4587}$, $13\mathrm{n}_{4589}$, $13\mathrm{n}_{4633}$, $13\mathrm{n}_{4634}$, $13\mathrm{n}_{4637}$, $13\mathrm{n}_{4638}$, $13\mathrm{n}_{4639}$, $13\mathrm{n}_{4640}$, $13\mathrm{n}_{4641}$, $13\mathrm{n}_{4642}$, $13\mathrm{n}_{4643}$, $13\mathrm{n}_{4644}$, $13\mathrm{n}_{4645}$}


\newtheorem{theorem}{Theorem}
\newtheorem{lemma}[theorem]{Lemma}
\newtheorem{corollary}[theorem]{Corollary}
\theoremstyle{definition}
\newtheorem{definition}[theorem]{Definition}
\theoremstyle{remark}
\newtheorem{remark}[theorem]{Remark}

\newcommand{\F}{\mathbb{F}}
\DeclareMathOperator{\rk}{rank}
\DeclareMathOperator{\Sp}{Sp}
\DeclareMathOperator{\GL}{GL}
\renewcommand{\O}{\mathrm{O}}

\title{Transvections and the bridge index of thirteen-crossing knots}
\author{Deep Bhattacharjee}
\email{itsdeep@live.com}
\date{October 20, 2026}


\begin{document}
\begin{abstract}
KnotInfo lists the bridge index of \ntotal{} prime knots with thirteen crossings as
either $3$ or $4$. We prove that \nclosed{} of them have bridge index $4$ and meridional
rank $4$, so the meridional rank conjecture holds for these knots. The lower bound comes
from representations of the knot group in $\GL_4(\F_q)$ that send a meridian to a
transvection or a reflection and fix no nonzero vector. Each such representation is
given explicitly and checked by a short program.
\end{abstract}
\maketitle

\section{Introduction}

Let $K\subset S^3$ be a knot, $G=\pi_1(S^3\setminus K)$ its group and $m\in G$ a meridian.
The \emph{bridge index} $b(K)$ is the least number of local maxima of a height function on
a knot isotopic to $K$. The \emph{meridional rank} $\mu(K)$ is the least number of
meridians, that is conjugates of $m^{\pm1}$, that generate $G$. The overpasses of a
$b$-bridge presentation give $b$ meridians generating $G$, so
\[
\mu(K)\le b(K).
\]
Cappell and Shaneson asked whether equality always holds
\cite[Problem~1.11]{Kirby}; this is the meridional rank conjecture.

Two lines of work make the conjecture checkable for a given knot. Blair, Kjuchukova,
Velazquez and Villanueva defined the Wirtinger number $\omega(D)$ of a diagram $D$ and
proved that its minimum over the diagrams of $K$ is $b(K)$ \cite{BKVV}. In particular
\begin{equation}\label{eq:upper}
b(K)\le\omega(D)\quad\text{for every diagram $D$ of $K$.}
\end{equation}
Blair, Kjuchukova and Morrison bounded $\mu(K)$ from below by surjections of $G$ onto
Coxeter groups sending meridians to reflections, and so computed $b(K)=\mu(K)$ for
595{,}515 knots up to sixteen crossings \cite{BKM}. KnotInfo \cite{KnotInfo} records the
bridge index of thirteen-crossing knots from these bounds and from the plain sphere
number \cite{BKP}. In the version 2026.10.5 of the database, \ntotal{} knots with thirteen
crossings have the entry $[3,4]$: the bounds $3\le b(K)\le4$ are known and the value is
not.

\begin{theorem}\label{thm:main}
For each of the \nclosed{} knots $K$ listed in Appendix~\ref{app:closed},
\[
b(K)=\mu(K)=4 .
\]
In particular the meridional rank conjecture holds for these knots.
\end{theorem}

The upper bound is \eqref{eq:upper} applied to the KnotInfo diagram of $K$, whose Wirtinger
number is $4$. The lower bound is Lemma~\ref{lem:fixed} below, applied to a representation
$\rho\colon G\to\GL_4(\F_q)$ with $q\in\{3,5\}$. Each $\rho$ sends the meridian to a
transvection of $\Sp_4(\F_q)$ or to a reflection of an orthogonal group $\O_4(\F_5)$. Reflections
in Coxeter groups are pseudo-reflections too; the transvections, of order $3$ and $5$, are
what settles most of these knots.
Table~\ref{tab:groups} gives the counts.

\begin{table}[ht]
\centering
\begin{tabular}{llr}
\toprule
group & image of a meridian & knots \\
\midrule
\grouptable
\bottomrule
\end{tabular}
\caption{The certificates of Theorem~\ref{thm:main}, by the group used.}\label{tab:groups}
\end{table}

\section{A lower bound from pseudo-reflections}

A \emph{pseudo-reflection} of a vector space $V$ is an element $g\in\GL(V)$ with
$\rk(g-1)=1$. It fixes the hyperplane $\ker(g-1)$ pointwise.

\begin{lemma}\label{lem:fixed}
Let $V$ be a vector space of dimension $n$ over a field and let $\rho\colon G\to\GL(V)$ be
a homomorphism such that $\rho(m)$ is a pseudo-reflection. Let $V^{G}$ be the subspace of
vectors fixed by $\rho(G)$. Then
\[
\mu(K)\ \ge\ n-\dim V^{G}.
\]
\end{lemma}

\begin{proof}
Let $m_1,\dots,m_k$ be meridians generating $G$, with $m_i=g_im^{\epsilon_i}g_i^{-1}$ and
$\epsilon_i=\pm1$. Since $\rho(m)^{-1}-1=-\rho(m)^{-1}(\rho(m)-1)$, each $\rho(m_i)$
is a pseudo-reflection, and $H_i=\ker(\rho(m_i)-1)$ is a hyperplane. The intersection
$H_1\cap\dots\cap H_k$ has dimension at least $n-k$ and is fixed by every $\rho(m_i)$,
hence by $\rho(G)$. So $n-k\le\dim V^G$.
\end{proof}

\begin{corollary}\label{cor:four}
If $\dim V=4$, $\rho(m)$ is a pseudo-reflection and $V^G=0$, then $b(K)\ge\mu(K)\ge4$.
\end{corollary}

In a diagram with arcs $x_1,\dots,x_r$, the Wirtinger presentation makes such a $\rho$ an
assignment of matrices $X_i=\rho(x_i)$ to the arcs that satisfies
$X_c=X_o^{\pm1}X_aX_o^{\mp1}$ at each crossing, where $x_o$ is the overarc and $x_a,x_c$
are the underarcs. Since the $X_i$ generate $\rho(G)$,
\[
V^G=\bigcap_{i=1}^{r}\ker(X_i-1).
\]
Both conditions are finite computations once the $X_i$ are written down.

\subsection*{Transvections}
Let $\omega$ be a nondegenerate alternating form on $V=\F_q^4$. For $v\ne0$ the
transvection $t_v(x)=x+\omega(x,v)\,v$ lies in $\Sp(V,\omega)$, $t_v-1$ has rank one,
$t_{-v}=t_v$, and $g\,t_v\,g^{-1}=t_{gv}$ for $g\in\Sp(V,\omega)$. A representation with
$X_i=t_{v_i}$ is therefore the same thing as a coloring of the arcs by vectors with
\[
v_c=\pm\,t_{v_o}^{\pm1}(v_a)
\]
at every crossing, and $\ker(X_i-1)=v_i^{\perp}$. So
$V^G=(\operatorname{span}\{v_i\})^{\perp}$, which is $0$ exactly when the $v_i$ span $V$.

\subsection*{Reflections}
Let $B$ be a nondegenerate symmetric bilinear form on $\F_q^4$, $q$ odd. For $B(v,v)\ne0$
the reflection $r_v(x)=x-2\frac{B(x,v)}{B(v,v)}\,v$ satisfies $g\,r_v\,g^{-1}=r_{gv}$ for
$g\in\O(B)$, and again $V^G=(\operatorname{span}\{v_i\})^{\perp}$.

\section{The certificates}\label{sec:cert}

For each knot of Appendix~\ref{app:closed} the file \texttt{data/certificates.json}
contains:
\begin{itemize}
\item the PD code of $K$ from KnotInfo, with edges numbered $1,\dots,2n$ along the
orientation and $[a,b,c,d]$ the edges at a crossing counterclockwise from the incoming
underedge $a$;
\item a prime $q$ and a $4\times4$ matrix over $\F_q$ for every edge;
\item four edges on distinct arcs, the seeds of a Wirtinger coloring.
\end{itemize}
The program \texttt{verify/verify.py} uses only the Python standard library. For every
knot it checks:
\begin{enumerate}
\item the matrices are constant along each arc, where an arc is a maximal run of edges
between two undercrossings;
\item at each crossing $[a,b,c,d]$ the relation $X_c=X_b^{-1}X_aX_b$ holds if the
overstrand runs from $b$ to $d$, and $X_c=X_bX_aX_b^{-1}$ if it runs from $d$ to $b$
(changing this convention at every crossing at once replaces $K$ by its mirror image, which
has the same $b$ and $\mu$);
\item each $X_i-1$ has rank $1$, and the matrices $X_i-1$ stacked together have rank $4$,
so that $V^G=0$;
\item starting from the four seed arcs, the coloring moves of \cite{BKVV} (if the overarc
and one underarc at a crossing are colored, color the other underarc) color every arc, so
$\omega(D)\le4$.
\end{enumerate}
By Corollary~\ref{cor:four} and \eqref{eq:upper}, items (1)--(4) prove $b(K)=\mu(K)=4$.
With the option \texttt{--knotinfo} the program also checks that each PD code is the one
in KnotInfo and that KnotInfo lists the bridge index as $[3,4]$. All \nclosed{}
certificates are checked in a few seconds. The program rejects altered
certificates: a changed matrix entry, a seed set on three arcs, one crossing with its
overstrand reversed, or identity matrices.

\begin{remark}
The certificates were found by a search, which plays no part in the proof. If the four
seed vectors do not span $V$, every arc vector lies in their span, so a certificate must
have seeds forming a basis. A coloring is then determined by the Gram matrix of the seeds:
six entries of an alternating matrix, or six off-diagonal entries of a symmetric matrix
with a fixed diagonal. The search runs over all such matrices and propagates the colors
from the seeds through the crossings. The Wirtinger relations are then checked at the
remaining crossings.
\end{remark}

\section{The knots that remain open}

For the other \nopen{} knots with entry $[3,4]$ (Appendix~\ref{app:open}) we found
neither bound. The exhaustive searches described above found no representation satisfying
Corollary~\ref{cor:four} with meridians sent to transvections of $\Sp_4(\F_q)$ for
$q\in\{3,4,5,7,8,9,11\}$, to reflections of $\O_4(\F_q)$ for $q\in\{3,5,7\}$ (both
isometry types), or to pseudo-reflections of $\GL_4(\F_2)$ and $\GL_4(\F_3)$ with
eigenvalue $1$ or $-1$. Random Reidemeister moves
(SnapPy \cite{SnapPy}) produced no diagram of Wirtinger number $3$; 300 diagrams were tried
for each knot. These knots stay open: a diagram of Wirtinger number $3$ would give
$b(K)=\mu(K)=3$, and a representation as in Corollary~\ref{cor:four} over another field,
or a Coxeter quotient of rank $4$, would give $b(K)=\mu(K)=4$.

\section*{Data availability}
The certificates, the verification program and the search programs are in the repository
\url{https://github.com/creelie/conjecture-to-choose}.

\begin{thebibliography}{9}
\bibitem{BKM} R.~Blair, A.~Kjuchukova and N.~Morrison, Coxeter quotients of knot groups
through 16 crossings, \emph{Experimental Mathematics} \textbf{34} (2025), no.~1, 27--37;
arXiv:2208.09032.
\bibitem{BKP} R.~Blair, A.~Kjuchukova and E.~Pfaff, The plain sphere number of a link,
arXiv:2504.10517.
\bibitem{BKVV} R.~Blair, A.~Kjuchukova, R.~Velazquez and P.~Villanueva, Wirtinger
systems of generators of knot groups, \emph{Comm. Anal. Geom.} \textbf{28} (2020), no.~2,
243--262; arXiv:1705.03108.
\bibitem{SnapPy} M.~Culler, N.~M.~Dunfield, M.~Goerner and J.~R.~Weeks, SnapPy, a computer
program for studying the geometry and topology of 3-manifolds, version 3.3.2,
\url{https://snappy.computop.org}.
\bibitem{Kirby} R.~Kirby (ed.), Problems in low-dimensional topology, in \emph{Geometric
Topology (Athens, GA, 1993)}, AMS/IP Stud. Adv. Math. \textbf{2.2}, Amer. Math. Soc.,
1997, 35--473.
\bibitem{KnotInfo} C.~Livingston and A.~H.~Moore, KnotInfo: table of knot invariants,
\url{https://knotinfo.org}, version 2026.10.5 (the \texttt{database\_knotinfo} package).
\end{thebibliography}

\appendix
\section{The knots of Theorem~\ref{thm:main}}\label{app:closed}
{\small\raggedright
\closedlist.\par}

\section{The knots with entry \texorpdfstring{$[3,4]$}{[3,4]} that remain open}\label{app:open}
{\small\raggedright
\openlist.\par}

\end{document}
